\documentclass[11pt,a4paper]{article}
\usepackage{curso}
\begin{document}
\portada{Lección}{8.2}{Matrices de curvatura}

\noindent\textbf{Unidad 8.}\quad\textbf{Objetivos:} Fisher empírico y Gauss--Newton generalizado (GGN); sistemas amortiguados; NGD-CG y NGD-Woodbury.
\textbf{Prerrequisitos:} 8.1; rango e identidad de Woodbury.

\section{Dos curvaturas}
Fisher empírico: $\bar F=JJ^\top=\frac1b\sum_i\nabla\log p_i\nabla\log p_i^\top$. GGN: $G=J_fH_LJ_f^\top$, con $J_f$ el jacobiano de la salida $f_\theta(x)$
y $H_L$ el hessiano de la pérdida respecto de la salida. Para pérdidas de la familia exponencial, $G$ coincide con el Fisher verdadero; $\bar F$ puede diferir mucho.
\begin{ejemplo}
Regresión logística, $N=200$, $b=32$, pesos no nulos \evS: $\|\bar F-G\|/\|G\|=232\,\%$. (En $w=0$ coinciden por casualidad: $(p-y)^2=p(1-p)=\frac14$.)
\end{ejemplo}

\section{Rango bajo y amortiguamiento}
$\mathrm{rank}\,G\le b\,d_z\ll N$ (en el ejemplo, rango 32): $G$ es singular y se resuelve $(G+\lambda I)\tilde g=\nabla\ell$.

\section{NGD-CG}
El gradiente conjugado usa solo productos $Gv=J_f(H_L(J_f^\top v))$, en $O(bN)$ con productos jacobiano--vector, sin formar $G$. Como $G+\lambda I$ tiene a lo sumo
$b\,d_z+1$ valores propios distintos, CG converge en a lo sumo $b\,d_z+1$ iteraciones en aritmética exacta: en el ejemplo, 21--23 iteraciones para
$\lambda=10^{-1},10^{-2},10^{-3}$ ($\kappa$ hasta $7{,}4\times10^3$). Las iteraciones son secuenciales.

\section{NGD-Woodbury}
Con $U=J_f$, $V=H_LJ_f^\top$: $(UV+\lambda I)^{-1}=\lambda^{-1}I-\lambda^{-2}U(I+\lambda^{-1}VU)^{-1}V$; solo se invierte una matriz $b\,d_z\times b\,d_z$. En el ejemplo
coincide con la solución directa hasta $10^{-12}$. Costo $O(b^2d_z^2N+b^3d_z^3)$: intratable si $d_z$ es un vocabulario ($\sim10^4$).

\section{Comparación (Donatella et al.)}
\begin{center}\small\begin{tabular}{lll}\toprule
optimizador & tiempo por iteración & memoria\\\midrule
SGD / Adam & $O(bN)$ & $O(N)$\\
NGD & $O(N^3+bd_zN^2)$ & $O(N^2)$\\
NGD-CG & $O(cbN)$ & $O(N)$\\
NGD-Woodbury & $O(b^2d_z^2N+b^3d_z^3)$ & $O(bd_zN+b^2d_z^2)$\\
NGD termodinámico & $O(bd_zN+t)$ & $O(bd_zN+b^2d_z^2)$\\\bottomrule\end{tabular}\end{center}

\begin{aplicacion}
Preferir GGN; Woodbury si $b\,d_z$ es pequeño; CG truncado si $N$ es grande. Laboratorio L11: NGD-CG y NGD-Woodbury en MNIST.
\end{aplicacion}

\begin{error}
\begin{itemize}[nosep]
\item Usar el Fisher empírico como el verdadero. \quad Formar $G$ cuando basta $Gv$.
\item Llevar CG a convergencia total. \quad Woodbury con $d_z$ = vocabulario.
\end{itemize}
\end{error}

\begin{thebibliography}{9}\small
\bibitem[40]{tngd} K. Donatella et al., Thermodynamic natural gradient descent, 2024.
\bibitem[37]{martens} J. Martens, New insights and perspectives on the natural gradient method, \emph{JMLR} 21 (2020).
\bibitem{schraudolph} N. N. Schraudolph, Fast curvature matrix-vector products for second-order gradient descent, \emph{Neural Computation} 14, 1723 (2002).
\end{thebibliography}
\end{document}
